\section{Cái bài ViT chứng minh, và cái nó không}
\label{sec:ch11-con-lai-gi}

\index{Vision Transformer!cái còn lại}%
ViT không chứng minh CNN đã lỗi thời. Nó làm một việc khó hơn: tách hai câu
người ta hay dán vào nhau.

Câu thứ nhất: attention toàn cục có đủ sức biểu diễn cho ảnh. Chương 10 đã có
một phiên bản hẹp của câu này: self-attention biểu diễn được convolution nếu
ta đặt đủ head đúng cách. ViT cho thấy, với JFT-300M, người ta không cần tự đặt
cách ấy bằng tay ở mỗi layer để đạt kết quả mạnh.

Câu thứ hai: vì thế cấu trúc ảnh không còn đáng giá. Bảng CIFAR-10 của mục
\ref{sec:ch11-cung-recipe} không cho phép nói vậy. CNN 66,570 tham số đạt
0.7036 ở 50,000 ảnh; ViT 66,095 tham số đạt 0.4576. ViT không phải MLP, nhưng
nó không thay CNN dưới recipe này. Cả hai câu cùng đúng vì chúng hỏi hai điều
khác: cái gì \emph{có thể học}, và cái gì \emph{được cho sẵn}.

Chỗ dễ lẫn cuối cùng là số trong bài với số tôi đo. Bài báo nói ViT-B/16 từ
77.91 lên 84.15 khi đổi pretraining data. Tôi không đo JFT, không dựng
ViT-B/16, và cũng không có ImageNet-1k trong companion repo. Tôi đo cái đối
nghịch: ViT nhỏ, CIFAR-10, CPU, không pretraining. Hai bảng ở hai đầu cùng một
đường, không phải hai phép tái lập của nhau.

Vậy chương này đóng được một câu hẹp: nếu bỏ tính cục bộ khỏi kiến trúc, dữ
liệu nhỏ và recipe ngắn không tự bù lại. Nó mở hai câu khó hơn.

Một là patch nên nhỏ tới đâu. Patch nhỏ giữ chi tiết tốt hơn và đẩy chi phí
lên bậc hai; patch lớn rẻ hơn nhưng đã nén ảnh trước attention. Hai là khi
attention đã phải nhìn 577 token hay hàng chục nghìn token, cái gì thực sự
giới hạn nó: số phép tính hay dữ liệu đi qua bộ nhớ? Chương sau không đổi bài
toán bằng một kernel mới. Nó đổi thứ attention phải mang qua bộ nhớ.
